\subsection{DET-11: an exact resolvent bridge for the scalar prime-sum endpoint}
There is a stronger affirmative comparison than the common residue denominator. It concerns the scalar endpoint itself, rather than an identification of the upstream finite matrices. Let $G_\zeta$ be exactly the function in \fid{DET-09}, and put
\[
 \Phi_\zeta(s)=\int_0^\infty G_\zeta(e^u)e^{-su}\,du,
 \qquad v=s(1-s),\qquad C=B(I+B/4)^{-1}.
\]
Write $\psi=\Gamma'/\Gamma$ and
\[
 A_\infty(s)=\frac1s+\frac1{s-1}-\frac12\log\pi+
 \frac12\psi(s/2).
\]
The following identities are unconditional, with the zero-based $B$ of \fid{B-01}.
\begin{proposition}
The trace-class determinant $D_C(v)=\det(I-vC)$ satisfies, for every $s\in\C$,
\begin{equation}\label{eq:Cdetcompletion}
 D_C(s(1-s))=2\xi(s).
\end{equation}
Where $\xi(s)\ne0$,
\begin{equation}\label{eq:Cresolvent}
 \frac{\xi'(s)}{\xi(s)}=(2s-1)\Tr\bigl(C(I-s(1-s)C)^{-1}\bigr).
\end{equation}
In particular, for $\Re s>1$ the actual scalar transform is
\begin{align}\label{eq:Phiresolvent}
 \Phi_\zeta(s)
 &=\frac{A_\infty(s)-(2s-1)\Tr\bigl(C(I-vC)^{-1}\bigr)}{s(s-1)}\notag\\
 &\hspace{7mm}-\frac1{(s-1)^2}+\frac{1+\gamma_E}{s-1}.
\end{align}
Thus its nontrivial-zero contribution is an ordinary trace of a resolvent of a fixed trace-class operator after explicit completion terms. This representation uses the zeros as input.
\end{proposition}
\begin{proof}
The diagonal entries of $I+B/4$ are nonzero: a zero would give $\gamma^2=-1/4$, hence a zero of $\xi$ at $0$ or $1$, whereas $\xi(0)=\xi(1)=1/2$. They tend to one, so the inverse is bounded. The entries of $C$ are therefore
\[
 c_j=(\gamma_j^2+1/4)^{-1}=[\rho_j(1-\rho_j)]^{-1},
 \qquad \sum_j|c_j|<\infty.
\]
Absolute locally uniform convergence of the products gives
\[
 D_C(v)=\frac{\prod_j(1-(v-1/4)\gamma_j^{-2})}
 {\prod_j(1+\gamma_j^{-2}/4)}
 =\frac{G(v-1/4)}{G(-1/4)}.
\]
Since $v-1/4=-(s-1/2)^2$, evenness of $\Xi$ gives
$G(v-1/4)=\xi(s)/\xi(1/2)$ and
$G(-1/4)=\xi(0)/\xi(1/2)$. This proves \eqref{eq:Cdetcompletion}, including the normalization factor two and with no square-root branch choice.

On any compact set avoiding the zeros of $D_C$, the sum
$\sum_j c_j/(1-vc_j)$ converges uniformly and absolutely. Indeed, the tail denominators are uniformly bounded away from zero because $c_j\to0$, and the finitely many remaining denominators are nonzero there. Differentiating the locally convergent logarithmic products consequently gives
\[
 \frac{D_C'(v)}{D_C(v)}=-\sum_j\frac{c_j}{1-vc_j}
 =-\Tr\bigl(C(I-vC)^{-1}\bigr).
\]
The inverse is bounded by the same diagonal argument, and its product with $C$ is trace class. The chain rule and $v'(s)=1-2s$ prove \eqref{eq:Cresolvent}. At $s=1/2$ no division by $2s-1$ is made; both sides of that identity vanish.

The source already computes the Mellin transform in
\path{Development/AutomorphicArithmetic/DedekindLFunctions/LandauTheorems/mellinGchiModOne.lean},
\texttt{mellin\_Gchi\_modOne}, lines 115--118. Its change from $y$ to $u=\log y$ is
\texttt{poleData\_change}, lines 114--135 of the adjacent \path{poleData.lean}.
For the modulus-one character its displayed formula is algebraically
\[
 \Phi_\zeta(s)=-\frac{\zeta'(s)/\zeta(s)}{s(s-1)}
 -\frac1{(s-1)^2}+\frac{1+\gamma_E}{s-1},\qquad \Re s>1.
\]
The modulus-one definition is \path{CharFamily.lean}, lines 15--22 in the parent directory; its ramp equals the sum in \path{Formalism/StatementOC.lean}, lines 11--13 for $y\ge1$, because terms beyond $\lfloor y\rfloor$ vanish and $\Lambda(0)=\Lambda(1)=0$.
Logarithmic differentiation of the defining completion gives
$\xi'/\xi=A_\infty+\zeta'/\zeta$. Substitute \eqref{eq:Cresolvent} to obtain \eqref{eq:Phiresolvent}. The convergence and transform statement are those of the cited source on $\Re s>1$; its formula agrees with direct integration of each ramp term, whose Mellin transform is $n^{-s}/[s(s-1)]$.
\end{proof}

This finding narrows the comparison: the source's scalar endpoint does have an exact Fredholm-resolvent representation, so an assertion that the connection stops at a similar denominator would be false. It does not identify any finite Gram, Capelli or period matrix with $C$ or $B$. It also does not make the source's scalar inequalities into a direct proof of selfadjointness of an arithmetic operator: the representation was built from the same zero data as \fid{B-01}. The passage from $C$ back to $B$ is the scalar spectral transform $B=C(I-C/4)^{-1}$, with bounded inverse because $c_j\ne4$ and $c_j\to0$. Geometrically, the determinant is still the analytic exterior-power construction of cote 1 and SGA 6, V, rather than a newly constructed cohomology over $\operatorname{Spec}\Z$.
